\chapter{Deret Fungsi dan Pertukaran Limit}

Sebuah pustaka numerik kecil menghitung $\ln(1+x)$ dengan deret
$\sum_{k\ge 1}(-1)^{k+1}x^{k}/k$. Penulisnya menguji fungsi itu pada $x=0.5$
dengan dua belas suku dan membandingkannya dengan pustaka standar: galatnya
$6.4\times 10^{-6}$. Uji itu lolos, dokumentasi menyebut deretnya akurat sampai
enam angka, dan kodenya masuk produksi. Pemakai pertama memanggilnya dengan
$x=0.99$. Dua belas suku memberi $0.6529204447$, sedangkan
$\ln 1.99=0.6881346387$. Galatnya $3.5\times10^{-2}$, sekitar lima ribu kali
lebih besar, pada fungsi yang sama dengan jumlah suku yang sama.

Kerusakan itu bukan kesalahan ketik. Untuk setiap $x$ di $(-1,1)$ deretnya
konvergen ke $\ln(1+x)$; yang berubah hanya berapa suku yang dibutuhkan: $16$
suku di $x=0.5$, $89$ di $x=0.9$, $721$ di $x=0.99$, $5248$ di $x=0.999$,
semuanya untuk galat di bawah $10^{-6}$. Tidak ada satu bilangan suku yang cukup
untuk semua $x$, sebab $x$ boleh sedekat apa pun ke $1$.

Bandingkan dengan potongan tertutup $[-r,r]$, $0<r<1$. Di sana satu anggaran
berlaku serentak untuk semua $x$: sisa deret berganti tanda dibatasi
$r^{N+1}/(N+1)$, dan batas itu tidak menyebut $x$. Untuk $r=0.5$ dan $N=12$
batasnya $9.4\times10^{-6}$, galat sebenarnya $6.4\times10^{-6}$. Inilah
perbedaan yang seluruh bab ini bicarakan.

Bab 5 mengurus deret bilangan, dan setiap $x$ memberi deret bilangan yang
berbeda. Bab 11 dan 13 mendefinisikan kekontinuan dan integral titik demi titik,
jadi keduanya tidak punya cara memindahkan limit melewati tanda integral tanpa
syarat tambahan. Alat yang belum ada adalah jarak antara dua fungsi; begitu
jarak itu ada, seluruh keluarga deret menjadi satu barisan di ruang metrik.

\begin{untukanak}
Bayangkan sebuah kelas dengan tiga puluh murid, dan guru menulis satu bilangan
di papan: seratus. Tugas setiap murid menebak bilangan itu, dan setelah tiap
tebakan guru mengumumkan apakah tebakan itu makin dekat. Murid pertama cepat
sekali: dalam lima tebakan ia sudah meleset kurang dari satu. Murid kedua
lambat: butuh dua puluh tebakan. Murid ketiga sangat lambat: butuh seratus.
Tetapi ketiganya, kalau diberi waktu cukup, akhirnya meleset kurang dari satu.

Sekarang guru mengganti pertanyaannya. Bukan "apakah setiap murid akhirnya
dekat", melainkan "adakah satu saat ketika semua murid sudah dekat pada saat
yang sama". Untuk kelas tadi jawabannya tidak: murid paling lambat selalu
tertinggal, dan selalu bisa ditambah murid yang lebih lambat lagi.

Pertanyaan yang sama muncul pada rumus yang dipakai komputer. Sebuah rumus
mungkin benar untuk setiap nilai yang dimasukkan, tetapi jumlah suku yang
dibutuhkan bisa membengkak tanpa batas begitu nilainya mendekati ujung. Kalau
program menganggap "benar untuk setiap nilai" sama dengan "satu rumus cukup
untuk semua nilai", program itu akan meleset jauh pada nilai yang paling
merepotkan, dan melesetnya tidak kelihatan pada uji yang dipilih penulisnya.

Di bab ini kita akan tahu mengapa dua pertanyaan tadi bukan soal kata-kata,
melainkan soal apakah dua operasi boleh ditukar urutannya: limit dengan
integral, limit dengan turunan, dan penjumlahan dengan keduanya.
\end{untukanak}

\section{Konvergensi titik demi titik dan konvergensi seragam}
\tingkat{dasar}

Ambil keluarga fungsi $f_1,f_2,\ldots$ yang semuanya terdefinisi pada satu
himpunan $E\subseteq\mathbb{R}$. Untuk satu titik $x$ yang kita pegang tetap,
nilai $f_1(x),f_2(x),\ldots$ hanyalah barisan bilangan, dan Bab 4 sudah
mengatakan apa artinya barisan itu konvergen. Karena itu definisi pertama hampir
menulis dirinya sendiri.

\begin{definisi}
Misalkan $f_n$ dan $f$ terdefinisi pada $E\subseteq\mathbb{R}$.
\begin{enumerate}
\item Kita katakan $f_n\to f$ \emph{titik demi titik}\index{konvergensi titik
      demi titik} pada $E$ bila untuk setiap $x\in E$ dan setiap
      $\varepsilon>0$ terdapat $N$ sehingga $n\ge N$ berakibat
      $|f_n(x)-f(x)|<\varepsilon$. Di sini $N$ boleh bergantung pada $x$ dan
      pada $\varepsilon$.
\item Kita katakan $f_n\to f$ \emph{seragam}\index{konvergensi seragam} pada
      $E$ bila untuk setiap $\varepsilon>0$ terdapat $N$ sehingga $n\ge N$
      berakibat $|f_n(x)-f(x)|<\varepsilon$ untuk \emph{semua} $x\in E$
      sekaligus. Di sini $N$ hanya bergantung pada $\varepsilon$.
\item Bila $h$ terbatas pada $E$, norma supremum\index{norma supremum} $h$
      adalah $\|h\|_\infty=\sup_{x\in E}|h(x)|$.
\end{enumerate}
\end{definisi}

\kamus{konvergensi seragam}{uniform convergence}{semua titik harus sudah dekat
pada saat yang sama, bukan satu per satu menurut jadwalnya masing-masing.}

Satu kuantor berpindah tempat, dan seluruh isi bab ini bergantung pada
perpindahan itu. Pada definisi pertama urutannya "untuk setiap titik, ada
ambang". Pada definisi kedua urutannya "ada satu ambang untuk semua titik".
Kalimat pertama benar untuk deret pada pembuka bab, kalimat kedua salah, dan
selisih keduanya adalah selisih antara $16$ suku dan $5248$ suku.

\begin{bacakeras}
Dua suara bergantian.

Suara pertama: ``Apakah deret ini sampai?''
Suara kedua: ``Sampai. Perlahan, tetapi sampai.''
Suara pertama: ``Sampai untuk titik ini?''
Suara kedua: ``Untuk titik ini. Untuk titik sebelah, lain kali, dengan tempo sendiri.''

Itulah janji yang pertama. Setiap titik boleh lambat, asal ia punya satu tempo
untuk dirinya sendiri. Titik yang jauh dari ujung sampai dalam belasan langkah;
titik yang menempel di ujung butuh ribuan.

Sekarang suara pertama bertanya dengan kata lain: ``Adakah satu tempo untuk
semua titik sekaligus?''

Suara kedua diam lebih lama. ``Di selang tertutup yang berhenti sebelum ujung,
ada. Di seluruh selang terbuka, tidak. Dan yang tidak ada itu bukan kegagalan
hitungan; ia hanya menuntut tempo yang bergerak bersama titiknya.''

Ulangi baris ini sampai terasa di lidah: semua titik boleh lambat, asal satu
tempo. Bila temponya harus berbeda untuk tiap titik, namanya bukan seragam.
\end{bacakeras}

\begin{contoh}
Misalkan $f_n(x)=x^n$ pada $E=[0,1]$. Untuk $0\le x<1$ berlaku $x^n\to0$,
sedangkan $x^n=1$ untuk semua $n$ bila $x=1$. Jadi $f_n$ konvergen titik demi
titik ke
\[
  f(x)=\begin{cases}0, & 0\le x<1,\\ 1, & x=1.\end{cases}
\]
Kekonvergenan itu tidak seragam, sebab $\sup_{x\in[0,1]}|x^n-f(x)|=1$ untuk
setiap $n$: di $x=1$ selisihnya tepat $1$, dan untuk $x$ yang cukup dekat ke $1$
dari kiri $x^n$ mendekati $1$ sementara $f(x)=0$.
\end{contoh}

Contoh itu memperlihatkan dua hal. Pertama, limit titik demi titik fungsi kontinu
tidak harus kontinu: setiap $x^n$ kontinu, limitnya melompat di $x=1$. Kedua,
keseragaman bergantung pada himpunannya, bukan hanya pada rumus fungsinya.

Ketergantungan itu bisa diukur. Pada $[0,r]$ dengan $r<1$ kekonvergenannya
seragam, sebab $\sup_{[0,r]}|x^n|=r^n\to0$. Yang menarik adalah bila batas
kanannya ikut bergerak bersama $n$. Pada $[0,1-1/n]$ supremumnya
$(1-1/n)^n\to1/e=0.3678794412$; jadi menyempit sebesar $1/n$ belum
menyelamatkan keseragaman. Pada $[0,1-1/\sqrt{n}]$ supremumnya
$(1-1/\sqrt n)^n$, yang runtuh seperti $e^{-\sqrt n}$: $2.2\times10^{-2}$ pada
$n=10$, $2.7\times10^{-5}$ pada $n=100$, $2.2\times10^{-44}$ pada $n=10^4$. Laju
penyempitan, bukan rumus fungsinya, yang menentukan.

Uji diagonal memberi keterangan tambahan yang lebih tajam. Ambil titik
$x_n=1-1/n$ yang mengejar titik putus. Maka $f_n(x_n)=(1-1/n)^n\to1/e$,
sedangkan $f(x_n)=0$ untuk semua $n$. Jadi bukan hanya supremumnya yang tidak
mengecil: nilai di titik-titik yang dipilih dengan tepat pun menolak menuju
limit titik demi titiknya. Pola ini dipakai berulang kali di bab ini. Untuk
menyangkal keseragaman, cari satu barisan titik yang membuat selisihnya tetap
besar.

Dua keluarga sederhana memisahkan kedua gejala itu. Yang pertama, $g_n$ adalah
segitiga bertinggi $1$ dengan alas $1/n$ dan puncak di $1-1/n$: puncaknya
bergerak ke kanan, tingginya tidak pernah turun, dan luasnya $1/(2n)$. Titik
$x<1$ hanya tersentuh puncak selama $1/(2(1-x))<n<3/(2(1-x))$, jadi di $x=0.9$
hanya sembilan nilai $n$ memberi nilai positif, di $x=0.99$ sembilan puluh
sembilan, dan setelah puncak terakhir lewat nilainya nol selamanya. Limitnya nol
di setiap titik, sedangkan $\|g_n\|_\infty=1$. Yang kedua, $h_n$ bertinggi $1/n$
dengan alas $1$ di tempat yang sama: puncaknya mengecil, bukan berpindah, dan
$\|h_n\|_\infty=1/n\to0$ (Gambar~\ref{fig:puncak}).

\begin{figure}[htbp]
\centering
\begin{tikzpicture}[x=5.6cm,y=3.1cm,font=\small]
  \draw[->,layergrey] (0,0)--(1.15,0) node[right,font=\scriptsize]{$x$};
  \draw[->,layergrey] (0,0)--(0,1.32) node[above,font=\scriptsize]{$y$};
  \node[below,font=\scriptsize,layergrey] at (1,0) {$1$};
  \foreach \n/\c in {16/layerindigo!15,8/layerindigo!30,4/layerindigo!50,2/layerindigo!75}{
    \fill[\c] ({1-1/\n-1/(2*\n)},0) -- ({1-1/\n},1) -- ({1-1/\n+1/(2*\n)},0) -- cycle;
  }
  \draw[dashed,layergrey] (0,1)--(1.06,1);
  \node[above,font=\scriptsize,layergrey] at (0.72,1) {tinggi $1$};
  \begin{scope}[xshift=7.0cm]
    \draw[->,layergrey] (0,0)--(1.15,0) node[right,font=\scriptsize]{$x$};
    \draw[->,layergrey] (0,0)--(0,1.32) node[above,font=\scriptsize]{$y$};
    \node[below,font=\scriptsize,layergrey] at (1,0) {$1$};
    \foreach \n/\c in {8/layerteal!18,4/layerteal!32,2/layerteal!52,1/layerteal!72}{
      \fill[\c] (0,0) -- (0.5,{1/\n}) -- (1,0) -- cycle;
    }
    \node[right,font=\scriptsize,layergrey] at (0.52,0.5) {$n=2$};
    \node[right,font=\scriptsize,layergrey] at (0.52,0.25) {$n=4$};
    \node[right,font=\scriptsize,layergrey] at (0.52,0.125) {$n=8$};
  \end{scope}
\end{tikzpicture}
\caption{Puncak yang bergerak dan puncak yang mengecil. Kiri: $g_n$ bertinggi
tetap $1$ dengan alas $1/n$ dan puncak di $1-1/n$; tingginya tidak turun,
sehingga $\|g_n\|_\infty=1$ walaupun $g_n(x)\to0$ di setiap titik. Kanan: $h_n$
bertinggi $1/n$ dengan alas $1$ di tempat yang sama; seluruh grafiknya masuk
pita setinggi $1/n$, sehingga konvergen seragam.}
\label{fig:puncak}
\end{figure}

\begin{lemma}
$f_n\to f$ seragam pada $E$ jika dan hanya jika $\|f_n-f\|_\infty\to0$.
\end{lemma}

\begin{proof}
Kekonvergenan seragam berarti: untuk setiap $\varepsilon>0$ ada $N$ sehingga
untuk $n\ge N$ berlaku $|f_n(x)-f(x)|\le\varepsilon$ pada semua $x\in E$,
yaitu $\sup_x|f_n(x)-f(x)|\le\varepsilon$. Itu persis pernyataan bahwa barisan
bilangan real $\|f_n-f\|_\infty$ menuju nol. Satu-satunya selisih teknis adalah
bahwa definisi memakai $<\varepsilon$ sedangkan norma memakai $\le\varepsilon$;
keduanya setara karena $\varepsilon$ sebarang.
\end{proof}

Lemma itu mengubah pertanyaan tentang fungsi menjadi pertanyaan tentang satu
barisan bilangan: memeriksa keseragaman cukup dengan satu supremum, menolaknya
cukup dengan menunjukkan supremum itu tidak menuju nol, seperti pada $x^n$.

Keseragaman punya gambar yang tepat, yaitu pita. Ambil $\varepsilon=0.1$. Bila
seluruh grafik $f_n$ untuk $n\ge N$ termuat di antara $f-\varepsilon$ dan
$f+\varepsilon$, maka satu $N$ itu bekerja untuk semua titik. Yang diperiksa
bukan satu titik, melainkan apakah ada bagian grafik yang keluar dari pita
selebar $0.2$. Pada $g_n$ pita itu selalu ditembus puncaknya; keadaan yang
diharapkan terlihat di Gambar~\ref{fig:pita}.

\begin{figure}[htbp]
\centering
\begin{tikzpicture}[x=2.55cm,y=4.0cm,font=\small]
  \draw[->,layergrey] (0,0)--(4.5,0) node[right,font=\scriptsize]{$x$};
  \draw[->,layergrey] (0,0)--(0,1.18) node[above,font=\scriptsize]{$y$};
  \fill[layergrey!12] plot[domain=0:4,samples=90] (\x,{0.55+0.22*sin(2*\x r)+0.1})
    -- plot[domain=4:0,samples=90] (\x,{0.55+0.22*sin(2*\x r)-0.1}) -- cycle;
  \draw[dashed,layergrey] plot[domain=0:4,samples=90] (\x,{0.55+0.22*sin(2*\x r)+0.1});
  \draw[dashed,layergrey] plot[domain=0:4,samples=90] (\x,{0.55+0.22*sin(2*\x r)-0.1});
  \draw[thick,accent] plot[domain=0:4,samples=90] (\x,{0.55+0.22*sin(2*\x r)});
  \draw[thick,layerindigo] plot[domain=0:4,samples=180]
    (\x,{0.55+0.22*sin(2*\x r)+0.06*cos(9*\x r)});
  \draw[<->,layergrey] (3.78,{0.55+0.22*sin(2*3.78 r)-0.1})
    -- (3.78,{0.55+0.22*sin(2*3.78 r)+0.1});
  \node[right,font=\scriptsize,layergrey] at (3.8,{0.55+0.22*sin(2*3.78 r)}) {$0.2$};
  \node[accent,font=\scriptsize,above left] at (0.42,{0.55+0.22*sin(2*0.42 r)}) {$f$};
  \node[layerindigo,font=\scriptsize,below right] at (0.95,{0.55+0.22*sin(2*0.95 r)+0.06*cos(9*0.95 r)}) {$f_n$};
\end{tikzpicture}
\caption{Pita seragam selebar $0.2$ di sekeliling $f$, yaitu di antara $f-0.1$
dan $f+0.1$. Grafik $f_n$ yang bergelombang masih termuat seluruhnya di dalam
pita, jadi $\|f_n-f\|_\infty\le0.1$ dan satu $N$ yang sama bekerja untuk semua
titik sekaligus.}
\label{fig:pita}
\end{figure}

\begin{tangan}
Ambil tali sepanjang dua meter, selotip, dan dua orang. Bentangkan tali lurus di
lantai, lalu tempel lima tanda berjarak sama. Orang pertama memegang ujung kiri
dan mengangkatnya $5$ cm saja. Lihat: hanya ujung itu yang naik, empat tanda
lain masih menempel di lantai. Itu kekonvergenan titik demi titik.

Sekarang kedua orang mengangkat kedua ujung bersama-sama dengan hitungan yang
sama: $5$ cm, lalu $3$ cm, lalu $1$ cm, lalu setengah sentimeter. Amati seluruh
panjang tali. Tidak ada titik yang tertinggal di lantai, dan jarak terjauh tali
ke lantai sama di semua titik, sebab keduanya diangkat serentak. Itu
kekonvergenan seragam.

Terakhir, orang pertama mengangkat ujungnya sambil berjalan ke tengah, sementara
orang kedua menurunkan ujungnya. Talinya tetap miring, dan selalu ada bagian
yang $5$ cm di atas lantai: setiap titik akhirnya turun, tetapi jarak terjauh
tidak pernah mengecil. Itulah $g_n$ pada Gambar~\ref{fig:puncak}.
\end{tangan}

\begin{lemma}[Kriteria Cauchy seragam]
Barisan $f_n$ konvergen seragam pada $E$ jika dan hanya jika untuk setiap
$\varepsilon>0$ ada $N$ sehingga $n,m\ge N$ berakibat
$|f_n(x)-f_m(x)|<\varepsilon$ untuk semua $x\in E$. Syarat terakhir ini disebut
kriteria Cauchy seragam\index{kriteria Cauchy seragam}.
\end{lemma}

\begin{proof}
$(\Rightarrow)$ Bila $f_n\to f$ seragam, ambil $N$ dari definisi kekonvergenan
untuk $\varepsilon/2$. Untuk $n,m\ge N$ dan setiap $x$,
\[
  |f_n(x)-f_m(x)|\le|f_n(x)-f(x)|+|f(x)-f_m(x)|<\varepsilon .
\]
$(\Leftarrow)$ Untuk setiap $x$ tetap, barisan $f_n(x)$ adalah barisan Cauchy
di $\mathbb{R}$, jadi konvergen ke suatu bilangan yang kita sebut $f(x)$. Ini
mendefinisikan sebuah fungsi $f$ titik demi titik. Untuk memperoleh
keseragamannya, tetapkan $n\ge N$ pada kondisi Cauchy dan biarkan
$m\to\infty$. Karena $f_m(x)\to f(x)$, diperoleh $|f_n(x)-f(x)|\le\varepsilon$
untuk semua $x\in E$ sekaligus, dan $N$ yang dipakai tidak bergantung pada $x$.
\end{proof}

Bukti arah kedua hanya memakai kelengkapan $\mathbb{R}$: ruang fungsi lengkap
terhadap norma supremum, dan kelengkapan itulah yang membuat deret fungsi bisa
dijumlahkan tanpa takut limitnya keluar dari kelas fungsi yang kita pakai.

\begin{catatan}
Kekonvergenan seragam adalah sifat pasangan: barisan dan himpunannya. Kalimat
"$x^n$ konvergen seragam" tidak lengkap tanpa menyebut pada himpunan mana. Pada
$[0,1]$ jawabannya tidak, pada $[0,0.99]$ ya: norma supremum galatnya
$0.99^n$, yaitu $0.37$ pada $n=100$ dan $4.3\times10^{-5}$ pada $n=1000$.
Karena itu setiap klaim keseragaman di bab ini menyebut himpunannya.
\end{catatan}

\section{Uji Weierstrass: keseragaman dari satu deret bilangan}
\tingkat{menengah}

Deret fungsi adalah deret yang sukunya fungsi. Jumlah parsialnya
$S_N=\sum_{n=1}^{N}f_n$ adalah fungsi, jadi deret fungsi tidak lain adalah
barisan fungsi yang khusus, dan semua definisi bagian sebelumnya berlaku begitu
kita mengganti $f_n$ dengan $S_N$.

\begin{definisi}
Deret $\sum_{n\ge1}f_n$ \emph{konvergen titik demi titik} pada $E$ bila untuk
setiap $x\in E$ deret bilangan $\sum_n f_n(x)$ konvergen. Deret itu
\emph{konvergen seragam}\index{deret fungsi} pada $E$ bila barisan jumlah
parsialnya konvergen seragam pada $E$.
\end{definisi}

Uji berikut adalah cara termurah memperoleh keseragaman, dan hampir semua
contoh di buku ini memakainya. Idenya: bila setiap suku dibatasi oleh bilangan
yang deretnya konvergen, maka ekor deret fungsi dibatasi oleh ekor deret
bilangan itu, dan ekor bilangan tidak tahu-menahu soal $x$.

\begin{teorema}[Uji-$M$ Weierstrass]\index{uji Weierstrass}
Misalkan $f_n$ terdefinisi pada $E$ dan ada bilangan $M_n\ge0$ dengan
$|f_n(x)|\le M_n$ untuk semua $x\in E$ dan semua $n$. Jika $\sum M_n$
konvergen, maka $\sum f_n$ konvergen seragam pada $E$. Bila $E$ adalah selang
tertutup terbatas dan setiap $f_n$ kontinu, maka jumlahnya kontinu pada $E$.
\end{teorema}

\begin{proof}
Tulis $S_N=\sum_{i=1}^{N}f_i$. Untuk $n>m$ dan setiap $x\in E$,
\[
  |S_n(x)-S_m(x)|\le\sum_{i=m+1}^{n}|f_i(x)|\le\sum_{i=m+1}^{n}M_i .
\]
Karena $\sum M_n$ konvergen, ekor $\sum_{i=m+1}^{n}M_i$ dapat dibuat lebih
kecil dari $\varepsilon$ asalkan $m,n$ cukup besar, dan pilihan itu tidak
menyebut $x$. Kriteria Cauchy seragam menutup bukti keseragamannya. Untuk
kekontinuan, setiap $S_N$ kontinu sebagai jumlah berhingga fungsi kontinu, dan
limit seragam fungsi kontinu kontinu; pernyataan terakhir itu dibuktikan di
bagian berikut.
\end{proof}

\begin{contoh}
Pada $E=[-1,1]$ deret $\sum_{n\ge1}x^n/n^2$ konvergen seragam, sebab
$|x^n/n^2|\le1/n^2$ dan $\sum1/n^2=\pi^2/6$ konvergen. Jumlahnya kontinu pada
$[-1,1]$, dan nilainya di $x=1$ adalah $1.6449340668$. Jumlah dua juta suku
memberi $1.6449335668$; selisihnya $5.0\times10^{-7}$, sebanding dengan batas
ekor $1/N=5\times10^{-7}$. Di $x=-1$ deretnya konvergen ke $-\pi^2/6$, sebab
suku-sukunya berganti tanda dengan nilai mutlak yang menurun menuju nol.
\end{contoh}

\begin{contoh}
Pada $E=\mathbb{R}$ deret $\sum_{n\ge0}2^{-n}\sin(2^n x)$ konvergen seragam,
sebab $|2^{-n}\sin(2^n x)|\le2^{-n}$ dan $\sum2^{-n}=1$. Setiap suku kontinu,
jadi jumlahnya kontinu di setiap titik. Yang mengejutkan, jumlah itu tidak
terdiferensialkan di titik mana pun: hasil klasik Weierstrass (1872), diperluas
Hardy sampai nisbah $1$, kita kutip tanpa bukti. Keseragaman menjamin
kekontinuan, tetapi tidak menjamin apa pun tentang turunan.
\end{contoh}

\begin{catatan}
Uji-$M$ memberi syarat cukup, bukan syarat perlu, dan kegagalannya bukan bukti
apa-apa. Deret $\sum_{n\ge1}(-1)^n(1-x)x^n$ pada $[0,1]$ konvergen seragam:
sisa setelah suku ke-$N$ dibatasi $(1-x)x^{N+1}$, yang maksimumnya pada $[0,1]$
adalah $\tfrac{1}{N+2}\left(\tfrac{N+1}{N+2}\right)^{N+1}\to0$, yaitu
$3.20\times10^{-2}$ pada $N=10$ dan $3.67\times10^{-4}$ pada $N=1000$. Tetapi
$\sup_{[0,1]}|f_n|=n^n/(n+1)^{n+1}$, sebanding dengan $1/(en)$, sehingga
$\sum_n\sup|f_n|$ divergen: $0.88$ sampai $N=10$ dan $2.53$ sampai $N=1000$.
Tidak ada deret pembanding yang konvergen, jadi uji-$M$ tidak terpakai walaupun
keseragamannya benar.
\end{catatan}

\begin{untukanak}
Ada dua cara mengerjakan sederet tugas. Cara pertama: selesaikan semua tugas
satu per satu, baru kemudian hitung totalnya. Cara kedua: hitung totalnya sambil
mengerjakan, lalu selesaikan sisanya. Untuk penjumlahan biasa kedua cara memberi
hasil yang sama, dan kita terbiasa menukarnya tanpa berpikir.

Tetapi ada tugas yang tidak suka ditukar. Bayangkan kamu mengukur panjang
bayangan sebuah tiang setiap jam selama satu hari, menghitung rata-ratanya, lalu
mencari jam ketika rata-rata itu paling panjang. Balik urutannya: cari dulu jam
terpanjang pada tiap hari, lalu rata-ratakan jam itu sepanjang bulan. Dua
hasilnya bisa berbeda jauh, bukan karena hitungannya salah, tetapi karena
"paling panjang" dan "rata-rata" memang tidak saling bertukar.

Hal yang sama terjadi pada rumus yang dipakai untuk menjumlahkan tak hingga
banyak suku, lalu mengintegralkan atau menurunkannya. Kadang urutannya boleh
ditukar dan hasilnya sama persis; kadang tidak, dan selisihnya besar. Di bab ini
kita akan tahu tanda bahayanya sebelum hasilnya meleset.
\end{untukanak}

\section{Kapan limit boleh ditukar}
\tingkat{menengah}

Tiga operasi yang biasa kita pakai pada fungsi adalah menghitung limit di
sebuah titik, mengintegralkan, dan menurunkan. Ketiganya sendiri adalah limit:
kekontinuan di $x_0$ berarti $\lim_{x\to x_0}f(x)=f(x_0)$, integral adalah
limit jumlah Riemann, turunan adalah limit hasil bagi selisih. Jadi memindahkan
limit barisan melewati salah satu operasi itu berarti menukar dua limit, dan
pertukaran itu tidak gratis: harganya satu hipotesis, yaitu
keseragaman\index{pertukaran limit}.

\begin{teorema}
Bila $f_n$ kontinu pada $E$ dan $f_n\to f$ seragam pada $E$, maka $f$ kontinu
pada $E$.
\end{teorema}

\begin{proof}
Ambil $x_0\in E$ dan $\varepsilon>0$. Karena kekonvergenannya seragam, ada $n$
dengan $\|f_n-f\|_\infty<\varepsilon/3$. Karena $f_n$ kontinu di $x_0$, ada
$\delta>0$ sehingga $|x-x_0|<\delta$ berakibat $|f_n(x)-f_n(x_0)|<\varepsilon/3$.
Maka untuk $x\in E$ dengan $|x-x_0|<\delta$,
\[
  |f(x)-f(x_0)|\le|f(x)-f_n(x)|+|f_n(x)-f_n(x_0)|+|f_n(x_0)-f(x_0)|
  <\varepsilon .
\]
\end{proof}

Perhatikan urutan logisnya: $n$ dipilih lebih dulu, baru $\delta$, dan
$\delta$ boleh bergantung pada $n$. Yang tidak boleh adalah memilih $\delta$
yang harus bekerja untuk semua $n$ sekaligus.

\begin{tantangan}
Deret $\sum_{n\ge0}\bigl(x^{2n}-x^{2n+1}\bigr)$ bersuku kontinu pada $[0,1]$
dan konvergen di setiap titiknya. Tentukan jumlahnya, lalu tunjukkan jumlah itu
tidak kontinu di $x=1$. Mengapa hasil ini tidak membantah teorema di atas?
Jawaban yang benar harus menyebut satu bilangan yang boleh bergantung pada
$\varepsilon$ tetapi tidak boleh bergantung pada $x$.

Petunjuk 1: jumlah parsialnya punya bentuk tertutup, dan bentuk itu memuat
$x^{2N+2}$.

Petunjuk 2: untuk $x<1$ jumlahnya $1/(1+x)$, sedangkan di $x=1$ setiap suku
bernilai nol; jadi nilainya melompat dari $1/2$ ke $0$.

Kalau mentok, kembali ke bagian uji diagonal dan tanyakan apa yang terjadi pada
$x_N=1-1/N$.
\end{tantangan}

\begin{teorema}[Integrasi suku demi suku]
Bila $f_n$ terintegralkan Riemann pada $[a,b]$ dan $f_n\to f$ seragam pada
$[a,b]$, maka $f$ terintegralkan pada $[a,b]$ dan
\[
  \int_a^b f=\lim_{n\to\infty}\int_a^b f_n .
\]
\end{teorema}

\begin{proof}
Untuk keterintegralannya, ambil $\varepsilon>0$ dan pilih $n$ dengan
$\|f_n-f\|_\infty<\varepsilon$. Karena $f_n$ terintegralkan, ada partisi
$\mathcal{P}$ dengan $U(\mathcal{P},f_n)-L(\mathcal{P},f_n)<\varepsilon(b-a)$.
Pada setiap subselang, $M_i(f)\le M_i(f_n)+\varepsilon$ dan
$m_i(f)\ge m_i(f_n)-\varepsilon$, sehingga
$U(\mathcal{P},f)\le U(\mathcal{P},f_n)+\varepsilon(b-a)$ dan
$L(\mathcal{P},f)\ge L(\mathcal{P},f_n)-\varepsilon(b-a)$. Ketiga ketaksamaan
itu memberi $U(\mathcal{P},f)-L(\mathcal{P},f)<3\varepsilon(b-a)$, dan kriteria
keterintegralan Bab 13 menutup argumennya. Untuk nilai integralnya, cukup
dicatat bahwa
\[
  \left|\int_a^b f_n-\int_a^b f\right|
  =\left|\int_a^b (f_n-f)\right|
  \le\int_a^b|f_n-f|\le(b-a)\,\|f_n-f\|_\infty ,
\]
dan ruas kanan menuju nol.
\end{proof}

\begin{akibat}[Integrasi deret]
Bila $\sum_{n\ge1}f_n$ konvergen seragam pada $[a,b]$ dan setiap $f_n$
terintegralkan di sana, maka
\[
  \int_a^b\sum_{n=1}^{\infty}f_n=\sum_{n=1}^{\infty}\int_a^b f_n .
\]
\end{akibat}

Akibat ini diterapkan pada jumlah parsial, dengan keseragaman biasanya
diperiksa lewat uji-$M$. Contohnya, deret $\sum_{n\ge1}x^n/n^2$ konvergen seragam
pada $[0,1]$, sehingga
\[
  \int_0^1\sum_{n=1}^{\infty}\frac{x^n}{n^2}\,dx
  =\sum_{n=1}^{\infty}\frac{1}{n^2(n+1)}=\frac{\pi^2}{6}-1=0.6449340668 .
\]
Angka terakhir juga diperoleh dengan menjumlahkan dua juta suku pertama deret di
ruas kanan, yang memberi $0.6449340668$ pada presisi mesin.

\begin{contoh}[Kegagalan tanpa keseragaman]
Misalkan $f_n(x)=nx(1-x^2)^n$ pada $[0,1]$. Untuk setiap $x\in(0,1]$ faktor
$(1-x^2)^n$ mengecil secara geometri, jadi $n(1-x^2)^n\to0$ dan $f_n(x)\to0$;
di $x=0$ semua $f_n$ bernilai $0$. Jadi limit titik demi titiknya fungsi nol.
Integralnya dapat dihitung tepat dengan substitusi $u=1-x^2$:
\[
  \int_0^1 nx(1-x^2)^n\,dx=\frac{n}{2}\int_0^1 u^n\,du=\frac{n}{2(n+1)} .
\]
Barisan itu menuju $1/2$, bukan $0$. Kuadratur numerik memberi
$0.4545454545$ untuk $n=10$ dan $0.4950495050$ untuk $n=100$, sama dengan
$n/(2(n+1))$ sampai digit terakhir. Jadi $\lim_n\int_0^1f_n=1/2$ sedangkan
$\int_0^1\lim_nf_n=0$.
\end{contoh}

Kegagalan itu punya bentuk geometris. Fungsi $f_n$ bernilai nol di kedua ujung
dan memuncak di $x=1/\sqrt{2n+1}$, dengan tinggi $1.3396685$ pada $n=10$ dan
$4.2834707$ pada $n=100$: puncaknya makin tinggi dan makin sempit, sehingga
luasnya tetap sekitar $1/2$ sementara tingginya di setiap titik menuju nol.

\begin{teorema}[Turunan suku demi suku]
Misalkan $f_n$ terdiferensialkan pada $[a,b]$, barisan $f_n'$ konvergen seragam
pada $[a,b]$, dan ada satu titik $x_0\in[a,b]$ tempat $f_n(x_0)$ konvergen.
Maka $f_n$ konvergen seragam ke suatu $f$ yang terdiferensialkan pada $[a,b]$,
dan
\[
  f'(x)=\lim_{n\to\infty}f_n'(x)
\]
untuk setiap $x\in[a,b]$.
\end{teorema}

\begin{proof}[Rangka bukti]
Intinya teorema nilai rata-rata pada selisih dua suku: untuk $m,n$ dan setiap
$x,y\in[a,b]$,
\[
  \bigl|(f_n-f_m)(x)-(f_n-f_m)(y)\bigr|\le\|f_n'-f_m'\|_\infty\,|x-y| .
\]
Karena $f_n'$ Cauchy seragam, ruas kanan dapat dibuat kecil serentak untuk semua
pasangan $x,y$ begitu $m,n$ besar. Bersama kekonvergenan di satu titik $x_0$,
ketaksamaan itu membuat $f_n$ Cauchy seragam, jadi $f_n\to f$ seragam. Untuk
turunan, bagi dengan $|x-y|$, biarkan $y\to x$, lalu tukar dua limit yang
muncul; ketaksamaan yang sama menjamin pertukaran itu sah. Rincian teknisnya ada
pada Rudin, Teorema 7.17.
\end{proof}

Syarat "ada satu titik tempat $f_n(x_0)$ konvergen" tidak dapat dibuang: $f_n$
dan $f_n+1$ punya barisan turunan yang sama, tetapi hanya satu di antaranya yang
konvergen di $x_0$. Sebaliknya, keseragaman $f_n$ saja juga tidak cukup untuk
menyimpulkan apa pun tentang turunan.

\begin{contoh}[Turunan yang tidak konvergen]
Misalkan $f_n(x)=\sin(nx)/n$ pada $[0,\pi]$. Karena $\|f_n\|_\infty=1/n$,
barisan itu konvergen seragam ke fungsi nol; nilainya $0.1$ pada $n=10$ dan
$0.01$ pada $n=100$. Tetapi $f_n'(x)=\cos(nx)$: di $x=0$ barisan itu bernilai
$1$ untuk semua $n$, sedangkan di $x=\pi/2$ nilainya $0,-1,0,1,0,-1,\ldots$
menurut paritas $n$, yang tidak konvergen. Jadi $f_n'$ tidak konvergen titik
demi titik di mana pun kecuali titik-titik khusus, apalagi seragam. Turunan
tidak memaafkan osilasi kecil berfrekuensi tinggi: $f_n$ mengecil menuju nol,
tetapi kemiringannya tidak.
\end{contoh}

Gambar~\ref{fig:turunan} memuat tiga kurva turunan itu pada $[0,\pi]$. Di
$x=\pi/2$ ketiganya bernilai $1$, $0$, dan $-1$: tiga nilai berbeda dari satu
barisan di satu titik, dan itu sudah cukup untuk mengatakan limitnya tidak ada.

\begin{figure}[htbp]
\centering
\begin{tikzpicture}[x=3.4cm,y=1.5cm,font=\small]
  \draw[->,layergrey] (0,0)--(3.55,0) node[right,font=\scriptsize]{$x$};
  \draw[->,layergrey] (0,-1.28)--(0,1.28) node[above,font=\scriptsize]{$y$};
  \draw[dashed,layergrey] (1.571,-1.25)--(1.571,1.25);
  \node[below,font=\scriptsize,layergrey] at (1.571,-1.25) {$\pi/2$};
  \draw[thick,layerindigo] plot[domain=0:3.1416,samples=160] (\x,{cos(4*\x r)});
  \draw[thick,layerteal] plot[domain=0:3.1416,samples=160] (\x,{cos(5*\x r)});
  \draw[thick,layerorange] plot[domain=0:3.1416,samples=160] (\x,{cos(6*\x r)});
  \fill[ink] (1.571,1) circle (1.5pt);
  \fill[ink] (1.571,0) circle (1.5pt) node[right,font=\scriptsize] {$0$};
  \fill[ink] (1.571,-1) circle (1.5pt) node[right,font=\scriptsize] {$-1$};
  \node[right,font=\scriptsize] at (1.585,1.06) {$1$};
\end{tikzpicture}
\caption{Turunan $f_n'=\cos(nx)$ pada $[0,\pi]$ untuk $n=4$ (indigo), $n=5$
(teal), dan $n=6$ (oranye). Di $x=\pi/2$ nilainya $1$, $0$, dan $-1$: tiga nilai
berbeda dari satu barisan di satu titik, jadi limit turunannya tidak ada,
walaupun $f_n=\sin(nx)/n$ konvergen seragam ke nol dengan amplitudo $1/4$,
$1/5$, dan $1/6$.}
\label{fig:turunan}
\end{figure}

\begin{catatan}
Dalam pembelajaran mesin, fungsi risiko sering ditulis sebagai harapan atas
distribusi data, $\mathbb{E}_{(x,y)\sim D}[\ell(f(x),y)]$, lalu diturunkan
terhadap parameter model. Bila distribusinya diskret berhingga, seperti pada
himpunan latih, harapan itu jumlah berhingga dan pertukaran turunan dengan
penjumlahan selalu sah tanpa syarat. Bila kontinu, harapan itu integral, dan
memindahkan turunan ke dalamnya menuntut syarat dominasi yang lebih halus
daripada keseragaman; alatnya adalah teorema kekonvergenan dominan dari teori
ukur, yang tidak dibahas di sini. "Turunkan di dalam integral" adalah operasi
berizin, bukan operasi yang selalu benar.
\end{catatan}

\section{Deret pangkat dan jari-jari konvergensi}
\tingkat{lanjut}

Deret pangkat adalah deret fungsi dengan bentuk paling kaku: suku ke-$n$ adalah
koefisien tetap dikali $x^n$. Kekakuan itu berbuah. Himpunan tempat deretnya
konvergen selalu sebuah selang yang simetris terhadap pusatnya, sehingga seluruh
informasi tentang di mana deret itu berguna termuat dalam satu bilangan.

\begin{definisi}
Deret berbentuk $\sum_{n\ge0}a_nx^n$ dengan $a_n\in\mathbb{R}$ disebut deret
pangkat\index{deret pangkat} berpusat di $0$. Bilangan
\[
  R=\sup\Bigl\{r\ge0:\ \sum_{n\ge0}a_nr^n\ \text{konvergen}\Bigr\}
\]
disebut jari-jari konvergensi\index{jari-jari konvergensi} deret itu.
\end{definisi}

\kamus{deret pangkat}{power series}{deret yang sukunya koefisien tetap dikali
pangkat naik dari satu variabel, seperti resep yang sama untuk semua nilai.}

\kamus{jari-jari konvergensi}{radius of convergence}{jarak terjauh dari pusat
tempat resep itu masih menghasilkan bilangan berhingga.}

\begin{teorema}
Misalkan $\sum a_nx^n$ deret pangkat dengan jari-jari $R$. Maka deret itu
konvergen mutlak untuk setiap $|x|<R$ dan divergen untuk setiap $|x|>R$.
Selain itu
\[
  \frac{1}{R}=\limsup_{n\to\infty}|a_n|^{1/n},
\]
dengan kesepakatan $1/0=\infty$ dan $1/\infty=0$.
\end{teorema}

\begin{proof}
Tetapkan $x\ne0$ dan tulis $|a_nx^n|^{1/n}=|x|\cdot|a_n|^{1/n}$. Menurut uji
akar Bab 5, deret $\sum a_nx^n$ konvergen mutlak bila
$\limsup_n|a_nx^n|^{1/n}<1$, yaitu bila $|x|<1/\limsup_n|a_n|^{1/n}$, dan
divergen bila limit superior itu lebih besar dari $1$, sebab sukunya tidak
menuju nol. Karena $x$ sebarang, batas antara kedua rejim itu tepat
$R=1/\limsup_n|a_n|^{1/n}$, dan himpunan kekonvergenannya termuat di $[-R,R]$.
\end{proof}

Pada $|x|=R$ uji akar tidak memberi keputusan, dan memang tidak ada keputusan
umum: $\sum x^n/n^2$ konvergen di kedua ujung, $\sum x^n/n$ hanya di $x=-1$, dan
$\sum x^n$ di keduanya divergen.

\begin{akibat}
Bila $a_n\ne0$ untuk $n$ cukup besar dan $|a_{n+1}/a_n|\to L$, maka $R=1/L$
dengan kesepakatan yang sama seperti di atas.
\end{akibat}

Alasannya adalah rantai ketaksamaan
\[
  \liminf|a_{n+1}/a_n|\le\liminf|a_n|^{1/n}\le\limsup|a_n|^{1/n}
  \le\limsup|a_{n+1}/a_n|,
\]
yang dibuktikan sebagai latihan di Bab 5. Bila nisbah punya limit, keempat
bilangan itu sama, jadi rumus nisbah dan rumus akar memberi $R$ yang sama. Dalam
hitungan tangan, rumus nisbah hampir selalu lebih cepat.

Sekarang bagian yang membuat deret pangkat berguna: ia boleh diturunkan dan
diintegralkan suku demi suku, dan hasilnya berlaku di seluruh selang
kekonvergenannya, bukan hanya di potongan kecilnya.

\begin{teorema}[Turunan dan integral suku demi suku]
Misalkan $\sum a_nx^n$ berjari-jari $R>0$. Maka untuk setiap $x$ dengan
$|x|<R$,
\[
  \frac{d}{dx}\sum_{n\ge0}a_nx^n=\sum_{n\ge1}n\,a_nx^{n-1},
  \qquad
  \int_0^x\sum_{n\ge0}a_nt^n\,dt=\sum_{n\ge0}\frac{a_n}{n+1}x^{n+1},
\]
dan kedua deret baru itu juga berjari-jari $R$.
\end{teorema}

\begin{proof}
Ambil $r$ dengan $0<r<R$ dan pilih $\rho$ dengan $r<\rho<R$. Deret
$\sum|a_n|\rho^n$ konvergen, jadi ada $C$ dengan $|a_n|\rho^n\le C$ untuk semua
$n$. Untuk $|x|\le r$,
\[
  |n\,a_nx^{n-1}|\le n\,\frac{C}{\rho}\left(\frac{r}{\rho}\right)^{n-1},
\]
dan deret $\sum_n n(r/\rho)^{n-1}$ konvergen sebab $r/\rho<1$. Uji-$M$ karena
itu memberi keseragaman deret turunan pada $[-r,r]$, dan teorema turunan suku
demi suku memberi kesamaan di sana. Setiap titik $|x|<R$ termuat pada suatu
$[-r,r]$ dengan $r<R$, jadi kesamaannya berlaku di seluruh $(-R,R)$. Untuk
integral, jalan yang sama dipakai pada $\sum|a_n|r^{n-1}$, yang juga konvergen.
Jari-jari tidak berubah karena koefisien deret turunan adalah $na_n$, dan
\[
  \limsup_n(n|a_n|)^{1/n}
  =\limsup_n|a_n|^{1/n}\cdot\lim_n n^{1/n}
  =\limsup_n|a_n|^{1/n},
\]
sebab $n^{1/n}\to1$. Argumen yang sama berlaku untuk koefisien $a_n/(n+1)$,
karena $(n+1)^{-1/n}\to1$.
\end{proof}

Bukti itu memakai seluruh isi bab: uji-$M$ untuk keseragaman, teorema turunan
untuk memindahkan limit, dan rumus jari-jari untuk memastikan deret baru tidak
menyempit. Keseragaman hanya diklaim pada $[-r,r]$, $r<R$; pada contoh pembuka,
$\ln(1+x)$ seragam pada setiap $[-r,r]$ dengan $r<1$ tetapi tidak pada $(-1,1)$.

\begin{contoh}
Untuk $|x|<1$, deret geometri memberi $\sum_{n\ge0}x^n=1/(1-x)$. Menurunkan
suku demi suku,
\[
  \sum_{n\ge1}n\,x^{n-1}=\frac{1}{(1-x)^2},
\]
dan pada $x=1/2$ ruas kanannya $4$. Jumlah dua ratus suku pertama deret di ruas
kiri memberi $4.000000$ pada presisi mesin. Mengintegralkan suku demi suku dari
$0$ ke $x$ memberi
\[
  \sum_{n\ge1}\frac{(-1)^{n+1}x^n}{n}=\ln(1+x),\qquad |x|<1,
\]
yaitu deret pada pembuka bab. Kedua hasil itu bukan kebetulan numerik: keduanya
konsekuensi teorema di atas, karena $1/2$ dan $x$ yang dipakai terletak di dalam
jari-jari.
\end{contoh}

Dengan teorema itu, pertanyaan pembuka bab terjawab. Deret $\ln(1+x)$ konvergen
seragam pada setiap $[-r,r]$ dengan $r<1$, sebab $|(-1)^{n+1}x^n/n|\le r^n/n$
dan $\sum r^n/n$ konvergen. Pada $(-1,1)$ keseragamannya gagal, dan harganya
terlihat pada jumlah suku: $16$ di $x=0.5$, $5248$ di $x=0.999$. Di titik ujung
$x=1$ deretnya tetap konvergen, sebab menjadi deret harmonik berganti tanda. Itu
bukan akibat jari-jari melainkan teorema Abel: bila $\sum a_nR^n$ konvergen,
maka $\sum a_nx^n$ konvergen seragam pada $[0,R]$. Kita kutip; buktinya ada di
Rudin, bab "Some Special Functions", Teorema 8.2.

Satu peringatan praktis: batas seragam dirancang untuk kasus terburuk. Pada
deret $\ln(1+x)$ dengan $N=12$ dan $r=0.5$, batas $r^{13}/13=9.4\times10^{-6}$
hanya tiga kali galat sebenarnya di $x=0.5$, tetapi di $x=0.1$ ia longgar sekitar
sepuluh ribu kali.

\begin{resep}
\begin{enumerate}
\item Tulis selisihnya, $M_N=\sup_{x\in E}|f_N(x)-f(x)|$, atau taksir dengan
      nilai di titik yang paling buruk.
\item Cari titik $x_N$ yang memaksimalkan selisih itu; biasanya puncak, atau
      titik yang mengejar ujung selang.
\item Hitung $M_N$ sebagai ungkapan yang hanya memuat $N$, lalu periksa
      $M_N\to0$.
\end{enumerate}
Peringatan urutan: cari $x_N$ lebih dulu untuk setiap $N$, baru hitung
supremumnya. Membalik urutan itu, yaitu memilih $N$ lalu mencari satu $x$ yang
mewakili semua titik, memberi kesimpulan seragam yang salah pada $g_n$ di
Gambar~\ref{fig:puncak}, yang puncaknya justru bergeser bersama $N$.
\end{resep}

\section{Contoh hitung}
\tingkat{menengah}

Tiga kelompok angka di bawah ini dihitung ulang di lab komputasi. Yang pertama
mengukur kegagalan keseragaman $x^n$ pada $[0,1]$.

\begin{center}
\small
\begin{tabular}{rrrrr}
\hline
$n$ & $\sup_{[0,1]}x^n$ & $(1-\tfrac1n)^n$ & $(1-\tfrac1{\sqrt n})^n$ & $0.9^{\,n}$\\
\hline
$10$    & $1$ & $0.3486784401$ & $2.234\times10^{-2}$ & $0.3486784401$\\
$100$   & $1$ & $0.3660323413$ & $2.656\times10^{-5}$ & $2.656\times10^{-5}$\\
$1000$  & $1$ & $0.3676954248$ & $1.108\times10^{-14}$ & $1.748\times10^{-46}$\\
$10000$ & $1$ & $0.3678610464$ & $2.249\times10^{-44}$ & $<10^{-450}$\\
\hline
\end{tabular}
\end{center}

Kolom kedua dan kelima adalah dua laju yang berbeda. Di titik tetap $x=0.9$
galatnya $0.9^n$, turun dengan faktor $0.9$ per suku: perlu $132$ suku untuk
mencapai $10^{-6}$, sedangkan di $x=0.5$ cukup $20$ suku dan di $x=0.99$ perlu
$1375$. Setiap titik punya laju eksponensialnya sendiri, dengan basis $x$. Kolom
kedua tidak bergerak: galat seragam pada $[0,1]$ tetap $1$ untuk semua $n$, jadi
laju penurunannya nol. Yang dibandingkan karena itu bukan dua bilangan,
melainkan sesuatu yang menuju nol dan sesuatu yang tidak.

Kelompok angka kedua adalah galat aproksimasi $e^{x}$ dengan deret Taylor di
$x=1$. Sisa deretnya dibatasi $\frac{e}{(N+1)!}$ untuk $x\in[0,1]$, sebab
turunan ke-$(N+1)$ dari $e^x$ adalah $e^x\le e$.

\begin{center}
\small
\begin{tabular}{rrrr}
\hline
$N$ & $S_N(1)=\sum_{k=0}^{N}\frac{1}{k!}$ & $e-S_N(1)$ & batas $\frac{e}{(N+1)!}$\\
\hline
$5$  & $2.716666666667$ & $1.615\times10^{-3}$ & $3.775\times10^{-3}$\\
$10$ & $2.718281801146$ & $2.731\times10^{-8}$ & $6.810\times10^{-8}$\\
$15$ & $2.718281828459$ & $5.063\times10^{-14}$ & $1.299\times10^{-13}$\\
$20$ & $2.718281828459$ & $0$ (di bawah presisi) & $5.320\times10^{-20}$\\
\hline
\end{tabular}
\end{center}

Batas itu selalu benar dan lebih besar daripada galat sebenarnya, dengan faktor
sekitar $2.3$ sampai $2.6$: sisanya kira-kira
$\frac{1}{(N+1)!}\left(1+\frac{1}{N+2}+\cdots\right)$, dan $e$ muncul karena
sukunya dijumlahkan di $x=1$. Pada $N=20$ galat sebenarnya tidak terukur lagi,
sebab $5.3\times10^{-20}$ jauh di bawah satu unit pembulatan terakhir pada
besaran $2.718$, yaitu sekitar $4.4\times10^{-16}$.

Kelompok angka ketiga adalah tiga jari-jari konvergensi, dihitung dengan dua
cara berbeda pada suku ke-$1000$.

\begin{center}
\small
\begin{tabular}{llrr}
\hline
deret & $R$ dari rumus & $|a_{1000}/a_{1001}|$ & $|a_{1000}|^{1/1000}$\\
\hline
$\sum x^n/n^2$          & $1$              & $1.0020010000$ & $0.9862794856$\\
$\sum \frac{n!}{n^n}x^n$ & $e=2.7182818285$ & $2.7169239322$ & $0.3694916635$\\
$\sum \frac{2^n}{n}x^n$  & $0.5$            & $0.5005000000$ & $1.9862320968$\\
\hline
\end{tabular}
\end{center}

Kolom ketiga menaksir $R$ dari nisbah dua suku berurutan, kolom keempat dari
kebalikan taksiran akar. Untuk deret kedua, $|a_n/a_{n+1}|=(1+1/n)^n$ naik ke
$e$; pada $n=1000$ nilainya $2.7169239$, meleset $0.0014$. Untuk deret pertama,
$(1+1/n)^2$ turun ke $1$ dari atas, jadi taksirannya $1.0020$, dan kekeliruan
$0.0020$ itu tanda kita baru di suku ke-$1000$. Kolom keempat lebih lambat menuju
limitnya, sebab $|a_n|^{1/n}$ mendekati $1$ dengan laju $|\ln a_n|/n$; pada
$n=1000$ selisihnya $0.0016$, sebanding dengan $\ln(2\pi n)/(2n)$ dikali $1/e$.

\kertas{Ringkasan angka. Galat seragam $x^n$ pada $[0,1]$ adalah $1$ untuk setiap
$n$, sementara galat di titik tetap $x$ adalah $x^n$: $132$ suku untuk $10^{-6}$
di $x=0.9$, $1375$ suku di $x=0.99$. Deret $e^x$ di $x=1$ mencapai galat
$2.7\times10^{-8}$ dengan $10$ suku, di bawah batas $6.8\times10^{-8}$.
Jari-jari tiga deret adalah $1$, $e$, dan $0.5$; integral
$\int_0^1 nx(1-x^2)^n dx=n/(2(n+1))$ bernilai $0.4545$ pada $n=10$, sementara
limit fungsinya nol.}

\section{Lab komputasi}
\tingkat{lanjut}

Kode berikut dijalankan apa adanya pada Python 3.14.4 dengan NumPy 2.5.2 dan
SciPy 1.18.0. Ia menghitung ulang ketiga kelompok angka di atas, ditambah
kuadratur untuk contoh tandingan integral.

\begin{lab}
\begin{verbatim}
import math
import numpy as np
from scipy.integrate import quad

print("n      sup[0,1]   (1-1/n)^n   (1-1/sqrt n)^n     0.9^n")
for n in (10, 100, 1000, 10000):
    a = (1 - 1/n)**n
    b = (1 - 1/math.sqrt(n))**n
    print(f"{n:6d} {1.0:9.6f} {a:11.6f} {b:17.3e} {0.9**n:10.3e}")

print()
print("suku ln(1+x) agar galat < 1e-6")
for x in (0.5, 0.9, 0.99, 0.999):
    N = 1
    while x**N / N > 1e-6:
        N += 1
    s = sum(((-1)**(k+1)) * x**k / k for k in range(1, N + 1))
    print(f"x={x:6.3f}  N={N:5d}  galat={abs(math.log(1+x) - s):.3e}")

print()
print("N      S_N              e - S_N        e/(N+1)!")
for N in (5, 10, 15, 20):
    s = math.fsum(1/math.factorial(k) for k in range(N + 1))
    b = math.e / math.factorial(N + 1)
    print(f"{N:3d} {s:.12f} {math.e - s:14.3e} {b:13.3e}")

print()
print("jari-jari dari nisbah suku ke-1000 dan ke-1001 (log-space)")
la = lambda n: -2 * math.log(n)
lb = lambda n: math.lgamma(n + 1) - n * math.log(n)
lc = lambda n: n * math.log(2) - math.log(n)
for nama, g in (("a", la), ("b", lb), ("c", lc)):
    R = math.exp(g(1000) - g(1001))
    H = math.exp(g(1000) / 1000)
    print(f"deret {nama}: R = {R!r}   |a_n|^(1/n) = {H!r}")
print("e =", repr(math.e))

print()
print("f_n(x) = n x (1-x^2)^n di [0,1]")
for n in (10, 100):
    val, _ = quad(lambda x, n=n: n * x * (1 - x*x)**n, 0, 1)
    print(f"n={n:4d}  integral={val!r}  n/(2(n+1))={n/(2*(n+1))!r}")
\end{verbatim}
\end{lab}

Keluaran yang diperoleh:

\begin{verbatim}
n      sup[0,1]   (1-1/n)^n   (1-1/sqrt n)^n     0.9^n
    10  1.000000    0.348678         2.234e-02  3.487e-01
   100  1.000000    0.366032         2.656e-05  2.656e-05
  1000  1.000000    0.367695         1.108e-14  1.748e-46
 10000  1.000000    0.367861         2.249e-44  0.000e+00

suku ln(1+x) agar galat < 1e-6
x= 0.500  N=   16  galat=3.049e-07
x= 0.900  N=   89  galat=4.478e-07
x= 0.990  N=  721  galat=4.915e-07
x= 0.999  N= 5248  galat=4.993e-07

N      S_N              e - S_N        e/(N+1)!
  5 2.716666666667      1.615e-03     3.775e-03
 10 2.718281801146      2.731e-08     6.810e-08
 15 2.718281828459      5.063e-14     1.299e-13
 20 2.718281828459      0.000e+00     5.320e-20

jari-jari dari nisbah suku ke-1000 dan ke-1001 (log-space)
deret a: R = 1.0020010000000001   |a_n|^(1/n) = 0.9862794856312105
deret b: R = 2.7169239322350776   |a_n|^(1/n) = 0.36949166347196
deret c: R = 0.5004999999999973   |a_n|^(1/n) = 1.9862320968418674
e = 2.718281828459045

f_n(x) = n x (1-x^2)^n di [0,1]
n=  10  integral=0.4545454545454544  n/(2(n+1))=0.45454545454545453
n= 100  integral=0.4950495049504951  n/(2(n+1))=0.49504950495049505
\end{verbatim}

Empat hal perlu dibaca dari keluaran ini. Pertama, kolom kedua tetap $1.000000$
untuk keempat nilai $n$, sementara kolom kelima sudah $3.5\times10^{-1}$ pada
$n=10$: kegagalan keseragaman terlihat sebagai angka. Kedua, baris terakhir kolom
kelima tercetak $0.000e+00$ bukan karena galatnya nol, melainkan karena
$0.9^{10000}=10^{-457.6}$ di bawah batas \texttt{float64} terkecil yang masih
normal: underflow.

Ketiga, kolom keempat memberi bukti numerik untuk dua limit yang sudah
dibuktikan: $(1-1/n)^n\to1/e=0.3678794$ dan $(1-1/\sqrt n)^n\to0$, yang runtuh
jauh lebih cepat.

Keempat, kuadratur SciPy memberi $0.4545454545454544$ untuk $n=10$, sedangkan
rumus $n/(2(n+1))$ memberi $0.45454545454545453$; selisih satu unit pembulatan
terakhir. Tanpa bentuk tertutup, kita tidak punya cara memastikan kuadratur tidak
mengukur hal lain.

Satu hal yang tidak terlihat di keluaran: semua batas di bab ini batas atas.
Baris $N$ yang dilaporkan kode dipilih dengan memeriksa $x^N/N$, jadi $N$ itu
dijamin cukup, bukan terkecil; galat sebenarnya $3.0\times10^{-7}$ sampai
$5.0\times10^{-7}$ di bawah $10^{-6}$.

\begin{untukdosen}
Urutan mengajar yang saya sarankan: mulai dari kegagalan pustaka $\ln(1+x)$ pada
pembuka, lalu definisi titik demi titik dan seragam dengan $x^n$ sebagai contoh
tunggal. Jangan buru-buru ke uji-$M$. Habiskan satu pertemuan penuh untuk
memastikan mahasiswa bisa mengatakan dengan tepat mengapa $x^n$ gagal seragam
dan mengapa $(1-1/\sqrt n)^n$ berhasil. Setelah itu uji-$M$, lalu tiga teorema
pertukaran, masing-masing dengan contoh tandingannya sendiri.

Dua miskonsepsi yang paling sering muncul. Pertama, mahasiswa mengira "setiap
$f_n$ kontinu dan $f_n\to f$" sudah cukup untuk menyimpulkan $f$ kontinu; mereka
lupa bahwa yang dibutuhkan keseragaman, dan $x^n$ adalah obatnya. Kedua,
mahasiswa mengira keseragaman turun-menurun: bila $f_n\to f$ seragam maka
$f_n'\to f'$. Contoh $\sin(nx)/n$ mematahkan itu dalam satu baris; sebaiknya
ditulis di papan sebelum teorema turunan muncul.

Aktivitas papan tulis yang berhasil: tulis empat barisan fungsi tanpa nama,
yaitu $x^n$, $x^n/n$, $nx(1-x^2)^n$, dan $\sin(nx)/n$, lalu minta kelas
menentukan untuk masing-masing apakah konvergen titik demi titik, seragam, dan
apakah integral serta turunannya bertukar. Jawabannya membentuk tabel $4\times4$
yang bisa diisi dalam dua puluh menit. Mahasiswa siap lanjut bila ia bisa
menjelaskan mengapa deret pangkat hanya seragam pada $[-r,r]$.
\end{untukdosen}

\section{Ringkasan}
\tingkat{dasar}

\begin{ringkasan}
\begin{itemize}
\item Kekonvergenan titik demi titik memberi setiap titik ambangnya sendiri;
      kekonvergenan seragam menuntut satu ambang untuk semua titik sekaligus,
      dan perbedaan itu setara dengan $\|f_n-f\|_\infty\to0$.
\item Uji-$M$ Weierstrass menyimpulkan keseragaman dari satu deret bilangan
      pembanding. Ia syarat cukup, bukan perlu; kegagalannya tidak membuktikan
      apa pun.
\item Limit seragam fungsi kontinu kontinu, dan integralnya boleh ditukar dengan
      limit. Turunan menuntut keseragaman barisan turunan ditambah kekonvergenan
      di satu titik.
\item Ketiga hipotesis itu tidak dapat dihapus: $x^n$, $nx(1-x^2)^n$, dan
      $\sin(nx)/n$ masing-masing mematahkan satu di antaranya.
\item Deret pangkat konvergen mutlak di dalam jari-jarinya dan divergen di luar;
      di dalam jari-jari ia boleh diturunkan dan diintegralkan suku demi suku
      tanpa mengubah jari-jarinya.
\item Keseragaman selalu lokal: $\ln(1+x)$ seragam pada setiap $[-r,r]$, $r<1$,
      tetapi tidak pada $(-1,1)$, dan harganya terlihat pada jumlah suku, dari
      $16$ di $x=0.5$ sampai $5248$ di $x=0.999$.
\end{itemize}
\end{ringkasan}

\section{Latihan}
\tingkat{menengah}

\begin{latihan}
\begin{enumerate}
\item Hitung dengan Python, untuk $n=10,100,1000$, nilai
      $\sup_{x\in[0,1]}|x^n|$ dan nilai $x^n$ di $x=0.9$. Jelaskan mengapa
      kolom pertama datar dan kolom kedua tidak.

\item Buktikan: bila $f_n\to f$ seragam pada $E$ dan setiap $f_n$ terbatas pada
      $E$, maka $f$ terbatas pada $E$ dan barisan $\|f_n\|_\infty$ terbatas.
      Tunjukkan dengan contoh bahwa kekonvergenan titik demi titik tidak cukup.

\item Misalkan $f_n(x)=x/(1+nx^2)$ pada $\mathbb{R}$. Tunjukkan bahwa
      $f_n\to0$ titik demi titik, tentukan $\|f_n\|_\infty$, dan buktikan bahwa
      kekonvergenannya seragam. Berapa $n$ terkecil sehingga
      $\|f_n\|_\infty<0.01$?

\item Nilai klaim ini: benar, salah, atau belum cukup bukti. "Karena $x^n\to0$
      untuk setiap $x\in[0,1)$, barisan $f_n(x)=x^n$ konvergen seragam pada
      $[0,1)$." Sebutkan tepat di mana penalaran itu patah, dan berikan satu
      barisan titik yang membuktikannya.

\item Hitung galat $|e-S_N(1)|$ untuk $N=5,10,15,20$ dan bandingkan dengan
      batas $e/(N+1)!$. Berapa $N$ terkecil agar batas itu di bawah
      $10^{-16}$? Apakah galat sebenarnya masih bisa diukur pada $N$ itu?

\item Buktikan bahwa deret $\sum_{n\ge0}x^n(1-x)$ konvergen titik demi titik
      pada $[0,1]$ tetapi tidak seragam, dengan menghitung jumlahnya lebih dulu
      dan memeriksa kekontinuannya.

\item Hitung jari-jari konvergensi $\sum x^n/n^2$, $\sum(n!/n^n)x^n$, dan
      $\sum 2^nx^n/n$ dari nisbah suku ke-$1000$ dan ke-$1001$, lalu bandingkan
      dengan nilai eksaknya. Deret mana yang taksirannya paling lambat, dan
      mengapa?

\item Buktikan bahwa bila $f_n\to f$ seragam pada $E$ dan $g$ terbatas pada $E$,
      maka $f_ng\to fg$ seragam pada $E$. Tunjukkan dengan contoh bahwa
      keterbatasan $g$ tidak dapat dihapus.

\item Untuk $f_n(x)=nx(1-x^2)^n$ pada $[0,1]$, hitung titik puncak, tinggi
      puncak, dan $\int_0^1f_n$ untuk $n=10,100$. Bandingkan tinggi puncak
      dengan $\sqrt{n/(2e)}$ dan jelaskan mengapa luasnya tidak menuju nol.

\item Turunkan $\sum_{n\ge0}x^n=1/(1-x)$ suku demi suku dan pakai hasilnya
      untuk menghitung $\sum_{n\ge1}n/2^n$. Berapa suku yang diperlukan agar
      jumlah parsialnya berada dalam $10^{-10}$ dari nilai eksaknya?
\end{enumerate}

\noindent\emph{Petunjuk untuk tiga butir tersulit.} Butir 4: sifat seragam
menuntut satu $N$ untuk semua $x$, jadi cukup tunjukkan bahwa untuk setiap $n$
ada $x$ dengan $x^n\ge\tfrac12$; ambil $x=2^{-1/n}$. Butir 6: jumlahkan deret
geometrinya lebih dulu; hasilnya fungsi yang berbeda dari limitnya di satu titik,
dan limit seragam fungsi kontinu harus kontinu. Butir 8: tulis
$|f_ng-fg|=|f_n-f|\,|g|$ dan pakai $\|g\|_\infty$; untuk contoh tandingannya
ambil $g(x)=1/x$ pada $(0,1)$ dan $f_n$ yang mengecil seragam tetapi tidak
cukup cepat.
\end{latihan}

\begin{untukpeneliti}
Yang terbukti dengan bukti lengkap di bab ini: kesetaraan keseragaman dengan
norma supremum, kriteria Cauchy seragam, uji-$M$ Weierstrass, ketiga teorema
pertukaran beserta contoh tandingannya, rumus jari-jari lewat limit superior, dan
ketetapan jari-jari pada penurunan serta pengintegralan suku demi suku. Dua hasil
dikutip tanpa bukti: ketidakdiferensialan fungsi Weierstrass (1872, diperluas
Hardy 1916) dan teorema Abel. Sumber primernya Rudin, \emph{Principles of
Mathematical Analysis}, bab "Sequences and Series of Functions" (7.1--7.17) dan
"Some Special Functions" (8.1--8.2); Abbott, \emph{Understanding Analysis}, bab
"Sequences and Series of Functions"; Stewart, \emph{Calculus}, bagian "Power
Series".

Yang sering dipopulerkan berlebihan adalah gagasan bahwa keseragaman
menghilangkan semua masalah. Ia menghilangkan tiga masalah tertentu, pada selang
tempat ia berlaku, dan tidak lebih. Pada deret Fourier keseragaman justru jarang,
dan pertukaran integral tetap sah karena alasan lain. Pada aproksimasi jaringan
saraf, yang dipakai adalah kekonvergenan pada himpunan hingga.

Pertanyaan terbuka: bila $f_n$ hanya diketahui lewat nilai pada grid berhingga,
seberapa besar kekeliruan taksiran $\|f_n\|_\infty$ oleh maksimum di grid? Tanpa
asumsi keteraturan tidak ada jawaban umum, dan untuk puncak sempit grid yang
tampak rapat bisa meleset seluruhnya. Batas keberlakuan: semua fungsi di sini
bernilai real pada himpunan bagian $\mathbb{R}$, dengan aritmetika 64-bit pada
Python 3.14.4, NumPy 2.5.2, dan SciPy 1.18.0.
\end{untukpeneliti}

\peta{Bab 14 pada Rudin, \emph{Principles of Mathematical Analysis}, bab
"Sequences and Series of Functions" (7.1--7.17) dan "Some Special Functions"
(8.1--8.2); pada Abbott, \emph{Understanding Analysis}, bab "Sequences and
Series of Functions"; pada Stewart, \emph{Calculus}, bagian "Power Series" dan
"Taylor and Maclaurin Series". Basisnya paruh kedua bab \emph{Integral, Deret
Fungsi, dan Pertukaran Limit} edisi pertama buku ini. Bab ini memakai definisi
limit dan kriteria Cauchy dari Bab 4, uji akar Bab 5, kekontinuan seragam Bab 11,
serta kriteria keterintegralan Bab 13, semuanya dirujuk, tidak diulang.}